\par\SourceLeaf{30}
{\rightskip=0pt plus 1em Das Besondere am Prinzip der virtuellen Geschwindigkeiten ist, dass es eine allgemeine Formel zur Lösung aller statischen Aufgaben liefert und damit alle anderen Prinzipien vertritt. Es weist seine Berechtigung dazu jedoch nicht so unmittelbar aus, dass es schon durch seine bloße Formulierung von sich aus plausibel erscheint.\par}

In \textit{dieser} Hinsicht scheint das Prinzip, das ich hier aufstellen werde, diesen Vorzug zu haben. Es hat aber auch einen zweiten Vorzug, indem es das Gesetz von Bewegung und Ruhe auf genau dieselbe Weise und in größter Allgemeinheit umfasst. Bei der allmählichen Entwicklung der Wissenschaft und beim Lernen des Einzelnen ist es durchaus angemessen, dass das Leichtere dem Schwereren, das Einfachere dem Verwickelteren und das Besondere dem Allgemeinen vorausgeht. Ist der Geist aber erst einmal auf dem höheren Standpunkt angelangt, verlangt er den umgekehrten Gang, bei dem die ganze Statik nur als ein ganz spezieller Fall der Mechanik erscheint. Selbst der oben erwähnte Geometer scheint darauf Wert zu legen. Er sieht es als Vorzug des Prinzips der kleinsten Wirkung an, dass es Gleichgewicht und Bewegung zugleich umfasst, wenn man das Gleichgewicht so ausdrückt, dass die lebendigen Kräfte in beiden Fällen Minimalwerte annehmen. Diese Bemerkung scheint aber eher geistreich als wahr, denn das Minimum tritt in beiden Fällen in ganz verschiedener Hinsicht auf.

Das neue Prinzip lautet nun wie folgt.

\textit{Ein System materieller Punkte, die auf beliebige Weise miteinander verbunden sind und deren Bewegungen zugleich beliebigen äußeren Beschränkungen unterliegen, bewegt sich in jedem Augenblick in größtmöglicher Übereinstimmung mit der freien Bewegung oder unter kleinstmöglichem Zwang. Als Maß für den Zwang, dem das ganze System in jedem Zeitabschnitt unterliegt, gilt dabei die Summe der Produkte aus dem Quadrat der Abweichung jedes Punkts von seiner freien Bewegung und seiner Masse.}

Es seien $m,m',m''$ usw. die Massen der Punkte, $a,a',a''$ usw. ihre Orte zur Zeit $t$ und $b,b',b''$ usw. die Orte, die sie nach dem unendlich
